\section{The local rational functional}
\label{s3-functional}

This section translates the rational-functional calculation in Fauzan,
Section~3, pp.~5--9, including (3.1)--(3.12).  The parameter $X$ remains
indeterminate: a valuation estimate must control every coefficient, including
the constant coefficient.  Write $v_p^G(A)$ for the minimum valuation of the
coefficients of $A\in\mathbb Q_p[X]$, with $v_p^G(0)=+\infty$.

\paragraph{\textnormal{[Derived] Pullback, reflection, and difference.}}
Let $L(x^d)=B_d$, with $B_1=-1/2$, and define
\begin{equation}
 \tau_9(P)=\frac{L(P^{(7)})}{8!},\qquad
 \tau_{9,X}\!\left(\frac1{x-r}\right)=H_{d(r)}^{(9)}-X,
 \qquad
 d(r)=\begin{cases}r,&r\geq0,\\-r-1,&r<0.\end{cases}
 \label{s3-tau}
\end{equation}
Extend by polynomial division and simple partial fractions.  The analogue
of (3.1) is
\begin{equation}
 \mu_{9,X}(R)=\tau_{9,X}\bigl(x^9R(-x^2)\bigr).
 \label{s3-pullback}
\end{equation}
Indeed, seven differentiations of $(-1)^e x^{2e+9}$ give
$(-1)^e B_{2e+2}\prod_{r=3}^{9}(2e+r)/8!$.  At a pole, division gives
\[
 \frac{x^9}{j^2-x^2}
 =-x^7-j^2x^5-j^4x^3-j^6x
  -\frac{j^8}{2}\left(\frac1{x-j}+\frac1{x+j}\right).
\]
Only $-x^7$ contributes to the polynomial functional, with value $-1/8$.
Using $H_j^{(9)}-H_{j-1}^{(9)}=j^{-9}$ therefore recovers
$j^8(X-H_j^{(9)})-1/8+1/(2j)$, including the easily missed constant term.

The Bernoulli identities $L(P(-1-x))=L(P(x))$ and
$L(P(x+1)-P(x))=P'(0)$ yield the analogues of (3.2)--(3.3):
\begin{align}
 \tau_{9,X}(g(-1-x))&=-\tau_{9,X}(g(x)),
 \label{s3-reflection}\\
 \tau_{9,X}(g(x+1)-g(x))&=\frac{g^{(8)}(0)}{8!}.
 \label{s3-difference}
\end{align}
The second identity requires $g$ to be regular at zero.  For integer poles,
reflection uses $d(-r-1)=d(r)$; the difference identity uses the same
harmonic increment, whose value for $g=1/(x-r)$ is $-r^{-9}$.

\subsection{Completion at a prime}

\paragraph{\textnormal{[Derived] Analytic and principal-part pieces.}}
Fix $p\geq11$.  Put
\[
 \kappa_d=\frac{d(d-1)\cdots(d-6)}{8!}B_{d-7}\quad(d\geq7),
 \qquad \kappa_0=\cdots=\kappa_6=0.
\]
Since $p\nmid8!$, von Staudt--Clausen gives $v_p(\kappa_d)\geq-1$.
Also $\kappa_d\in\mathbb Z_p$ for $d\leq p+5$, because then the relevant
Bernoulli index is at most $p-2$.  This is a sufficient degree cutoff;
no optimal cutoff is needed here.  Consequently
\[
 \tau_9\!\left(\sum_{d\geq0}f_dz^d\right)=\sum_{d\geq0}f_d\kappa_d,
 \qquad \|\tau_9\|\leq p,
 \qquad \mathcal A=\mathbb Q_p\langle z\rangle .
\]
Thus the analytic part can lose one power of $p$.

Let $T(z)=\prod_\nu(z-r_\nu)$, where the integer roots have pairwise
unit differences and $d(r_\nu)<p$.  Polynomial division and partial
fractions have norm at most one, since
$c_\nu=U(r_\nu)/T'(r_\nu)$ and $T'(r_\nu)$ is a unit.  They extend to
the direct sum
\[
 \mathcal B_T=\mathcal A\oplus\bigoplus_\nu
       \mathbb Q_p\frac1{z-r_\nu},\qquad
 \tau^{\mathrm{ext}}_{9,Y}\!\left(f+\sum_\nu\frac{c_\nu}{z-r_\nu}\right)
 =\tau_9(f)+\sum_\nu c_\nu(H_{d(r_\nu)}^{(9)}-Y).
\]
For $Y\in\mathbb Z_p[X]$, the principal-part values are integral.
Accordingly, the analogue of (3.4) is
\begin{equation}
 \tau^{\mathrm{ext}}_{9,Y}
 \!\left(\sum_{j\geq0}\frac{p^jU_j(z)}{T(z)}\right)
 \in\mathbb Z_p[X],\qquad
 U_j\in\mathbb Z_p[z],\quad\deg U_0\leq p+5.
 \label{s3-completion}
\end{equation}
The initial quotient has integral functional value by its degree; every
later term carries the extra $p$ needed to absorb the analytic loss.
Convergence follows in the maximum norm on $\mathcal B_T$.

For a far pole, $v_p(s)<0$, expand $1/(z-s)=-\sum_{d\geq0}z^ds^{-d-1}$.
Applying $\tau_9$ termwise gives the explicit analogue of (3.5):
\begin{equation}
 \tau_9^{\mathrm{an}}\!\left(\frac1{z-s}\right)
 =-\frac18\sum_{k\geq0}\binom{k+7}{7}B_k\,s^{-k-8}.
 \label{s3-far-pole}
\end{equation}
The series converges because $v_p(s^{-1})>0$ and the Bernoulli valuations
are bounded below.  The decomposition into analytic functions and the
specified principal parts is unique.  Hence expanding far factors first
and taking global partial fractions first produce the same value.

\subsection{Distribution}

\paragraph{\textnormal{[Derived] The new scale is $p^{-8}$.}}
Define the analogues of the constants in (3.6) by
\begin{equation}
 C_{p,9}=\sum_{a=1}^{p-1}\tau_9^{\mathrm{an}}
                  \!\left(\frac1{z+a/p}\right),
 \qquad Y=p^9X+C_{p,9}.
 \label{s3-distribution-constant}
\end{equation}
Here $C_{p,9}\in p^9\mathbb Z_p$.  In fact, the $k=0$ part of
\eqref{s3-far-pole} is $-p^8\sum_{a=1}^{p-1}a^{-8}/8$, and its residue
vanishes because $p-1\nmid8$ for $p\geq11$.  The $k=1$ part already
has valuation at least $9$; for $k\geq2$ the bound is $k+7\geq9$.
The condition $p\geq11$ also makes $8!$ a unit.  Smaller primes are
handled separately below.

The analogue of (3.7), for rational $g$ with simple integer poles, is
\begin{equation}
 \tau_{9,X}(g)=p^{-8}\sum_{a=0}^{p-1}
        \tau^{\mathrm{ext}}_{9,Y}\bigl(g(a+pz)\bigr).
 \label{s3-distribution}
\end{equation}
Bernoulli multiplication contributes $p^{-1}$ and the seven derivatives
contribute $p^{-7}$, proving the polynomial identity.  At $g=1/x$ the
right side is $p^{-9}(-Y+C_{p,9})=-X$, which explains the parameter choice.
The substitutions $a\mapsto p-1-a$, $z\mapsto-1-z$ preserve reflection.
Translation telescopes to
\[
 p^{-8}\tau^{\mathrm{ext}}_{9,Y}
       \bigl(g(p(z+1))-g(pz)\bigr)=\frac{g^{(8)}(0)}{8!}.
\]
The analytic identities extend by continuity.  Starting with the pole
at zero, translation determines positive integer poles and reflection
determines negative ones.  This establishes the whole rational identity,
including its polynomial contribution.

\subsection{Small primes and the integer-valued basis}

\paragraph{\textnormal{[Derived] A uniform, deliberately coarse bound.}}
For any prime $p$, an integer-valued polynomial $A$ of degree at most $d$
has integral coefficients in the basis $\binom{x}{k}$, given by forward
differences.  Vandermonde's identity gives the unchanged analogue of (3.8):
\[
 v_p(A(x)-A(y))\geq v_p(x-y)-\lfloor\log_p\max(1,d)\rfloor.
\]
Its $r$th derivative loses at most $r\lfloor\log_p\max(1,d)\rfloor$:
the coefficient of $z^r$ in $\binom{z}{j}$ is a sum of products of $r$
reciprocals of integers at most $j$.  Moreover
$L(\binom{x}{k})=(-1)^k/(k+1)$, from
$L((1+z)^x)=\log(1+z)/z$.  Seven derivatives followed by $L$ therefore
give the analogue of (3.9):
\begin{equation}
 P(\mathbb Z_p)\subset p^{-A}\mathbb Z_p,\quad\deg P\leq d
 \quad\Longrightarrow\quad
 v_p(\tau_9(P))\geq-A-8\lfloor\log_p(d+1)\rfloor-v_p(8!).
 \label{s3-small-polynomial}
\end{equation}

For integer-valued $A$, $\deg A\leq d$, set
\[
 g(x)=\frac{(K!)^2A(x)}{\prod_{-K\leq r\leq K,\ r\ne0}(x-r)},
 \qquad L_0=\lfloor\log_p(2K)\rfloor,
 \qquad M=\lfloor\log_p\max(2K,d+1)\rfloor.
\]
Its residue at $r$ is
$c_r=(-1)^{K-r}(K!)^2rA(r)/((K+r)!(K-r)!)$, of valuation at least
$-L_0$.  The root-counting and inverse-product estimates of source
Lemma~3.3 are independent of the zeta index.  Outside the balls
$r+p^{L_0+1}\mathbb Z_p$ they bound the rational terms by $-2L_0$;
inside a ball, subtracting its principal part and using the preceding
difference estimate bounds the polynomial quotient $P$ by $-2M$.
Applying \eqref{s3-small-polynomial} gives a loss at most $10M+v_p(8!)$.
The pole terms have the same bound, since
$v_p(H_{d(r)}^{(9)})\geq-9\lfloor\log_pK\rfloor$.
Thus the analogue of (3.10) is
\begin{equation}
 v_p^G(\tau_{9,X}(g))\geq-10M-v_p(8!).
 \label{s3-small-rational}
\end{equation}

Keep exactly the source basis and normalization:
\[
 q_0=1,\qquad q_i(t)=\frac{(-1)^i2tD_{i-1}(t)}{(2i)!}\ (i\geq1),
 \qquad f_i(t)=\left(\frac{D_N(t)}{(N!)^2}\right)^3q_i(t).
\]
Both identities
\[
 q_i(-x^2)=\binom{x+i}{2i}+\binom{x+i-1}{2i},\qquad
 \frac{D_N(-x^2)}{(N!)^2}=\binom{N-x}{N}\binom{N+x}{N}
\]
show that $f_i(-x^2)$ is integer-valued on $\mathbb Z_p$.
Squaring the determinant of this triangular basis change gives precisely
the prescribed $S_K$.  Hence the analogue of (3.11) is
\begin{equation}
 F_{9,K}(X)=\det\left[
 (K!)^2\mu_{9,X}\!\left(\frac{f_i(t)f_j(t)}{D_K(t)}\right)
 \right]_{0\leq i,j<h}.
 \label{s3-normalized-basis}
\end{equation}
The pullback numerator, up to the harmless sign $(-1)^K$, is
$x^9f_i(-x^2)f_j(-x^2)$.  Its degree is at most
$12N+4h+5=23K/5+5$.  For $K\in40\mathbb Z_{>0}$ this gives
$\max(2K,d+1)\leq5K$.  Entrywise use of
\eqref{s3-small-rational} and the determinant expansion finally yield
the analogue of (3.12):
\begin{equation}
 \boxed{\displaystyle
 v_p^G(F_{9,K})\geq
 -10h\lfloor\log_p(5K)\rfloor-hv_p(8!)
 \qquad\text{for every prime }p.}
 \label{s3-small-final}
\end{equation}
This accounts for small-prime denominators.  Whether the sharper
large-prime estimates leave enough real decay is a separate question.
